\scanpage{37}
\begingroup\makeatletter\tagsleft@true\makeatother
\uline{Cor} \deleted{D}ans le cas où $\operatorname{Card}\{\ \}=4$, \inserted{et $\Sigma_1^-\ne X,\Sigma_2^-\ne X$} les seules relations d'inclusions entre les ensembles $\Sigma_1^+,\Sigma_1^-,\Sigma_2^+,\Sigma_2^-$ sont les relations
% Slanted marginal reformulation.
\margin{Sous forme : Pour que les relations \unsure{strictes} $(*)$ soient les seules relations d'inclusion entre les ensembles $\Sigma_1^+,\Sigma_1^-,\Sigma_2^+,\Sigma_{\unsure{1|2}}^-$, il f et s. que ceux-ci soient \inserted{\unsure{mutlt.}} distincts et que l'on ait $\Sigma_1^-\ne X,\ \Sigma_2^-\ne X$.}
\[
\left\{\begin{aligned}
\Sigma_1^+&\subset\Sigma_2^-\\
\Sigma_2^+&\subset\Sigma_1^-
\end{aligned}\right.\qquad .
\tag*{($*$)}
\]
En effet, on n'a pas les inclusions inverses, par hypothèse. D'autre part, on n'\uline{a pas}
\[
\Sigma_1^+\subset\Sigma_1^-\quad\text{ou}\quad\Sigma_1^-\unsure{\subset|\in}\Sigma_1^+,\quad\text{car on aurait}
\]
$\deleted{\unsure{\Sigma_1^+,\Sigma_1^+|\Sigma_1^+,\Sigma_1^-}}\ \Sigma_1^-=X$, de $=$ on n'a pas $\Sigma_2^+\subset\unsure{\deleted{\vert}}\unsure{\Sigma|E}_{\unsure{2|1}}^-$ ni $\Sigma_2^-\subset\Sigma_2^+$, et on n'a pas $\Sigma_1^+\subset\Sigma_2^+$ ni l'inclusion inverse, puisque $\Sigma_1^+$ et $\unsure{\Sigma|\Sigma_2}^{\unsure{\top|+}}$ sont disjoints et $\ne\varnothing$. Il reste à voir que l'on n'a pas une inclusion $\Sigma_1^-\subset\Sigma_2^-$ ou inclusion inverse, on lit sur le diagramme $(14)$ que \deleted{\unsure{cela}} \inserted{la première \unsure{p.ex.|p. ex.}} signifierait que $\deleted{\Sigma_{12}\simeq\Sigma_1^-}$ \deleted{\unsure{d'où}} $\Sigma_2^-=X$, ce que l'on a exclu.

% The long rule is the manuscript's left-margin emphasis beside the definition.
\noindent\makebox[0pt][r]{\rule[-30ex]{.4pt}{31.5ex}\hspace{.55em}}\uline{Déf} On dit que $\unsure{\Pi|\pi}_1,\unsure{\Pi|R}_2$ sont \uline{\unsure{strictement|Strictement} non panachées} \inserted{si elles sont non panachées, et} si l\unsure{\deleted{es}}es quatre ensembles $\Sigma_1^+,\Sigma_1^-,\Sigma_2^+,\Sigma_2^-$ sont distincts, et. si \unsure{((|c} de plus $\Sigma_1^-\ne X,\Sigma_2^-\ne X$ -- i.e si les seules relations d'inclusion entre les ensembles $\Sigma_1^+,\Sigma_1^-,\Sigma_2^+,\Sigma_2^-$ sont $\Sigma_1^+\subset\Sigma_2^-,\ \Sigma_2^+\subset\Sigma_1^-$ -- ce qui signifie aussi que l'on a :
\[
\Sigma_1^-\ne X,\Sigma_2^-\ne X
\quad(\text{on dit donc que }\Pi_{\unsure{1|2}},\Pi_2\text{ sont \uline{propres}}),
\ E_1\ne\Sigma_{12},E_2\ne\Sigma_{12}.
\tag*{(21)}
\]
\endgroup
\clearpage
% archive sheet number: 36
